\chapter{Υπολογιστική και Προγραμματιστική Υλοποίηση του Προσομοιωτή}
\label{ch:computational-implementation}

\noindent\textbf{Σύνοψη Κεφαλαίου}.
Το κεφάλαιο αυτό δείχνει πώς περνάμε από τη διεπαφή του προσομοιωτή στον
αριθμητικό πυρήνα. Πρώτα παρουσιάζονται η τεχνολογική στοίβα και η βασική
αρχιτεκτονική της εφαρμογής, έπειτα η ροή χρήσης των tabs και τέλος οι κύριες
εκτελεστικές διαδρομές για τροχιές, διαγράμματα διακλάδωσης και φάσματα
Lyapunov. Οι πιο τεχνικές λεπτομέρειες για πολυπλοκότητα, μνήμη και
Numba/JIT δίνονται στο Παράρτημα.

\section{Σκοπός του κεφαλαίου}
\label{sec:ch2-purpose}

Το πρώτο κεφάλαιο ασχολήθηκε κυρίως με τις έννοιες και τα βασικά εργαλεία της
ανάλυσης. Στο παρόν κεφάλαιο το ενδιαφέρον μεταφέρεται στον τρόπο με τον οποίο
αυτές οι έννοιες υλοποιούνται σε λογισμικό.

Ο στόχος εδώ δεν είναι να παρουσιαστεί όλος ο κώδικας γραμμή προς γραμμή.
Στόχος είναι να γίνει σαφές:
\begin{itemize}
\item ποια είναι τα βασικά μέρη της εφαρμογής,
\item πώς ο χρήστης δίνει τις ρυθμίσεις ενός πειράματος,
\item πώς οι ρυθμίσεις αυτές φτάνουν στον αριθμητικό πυρήνα,
\item και ποια στοιχεία της τρέχουσας υλοποίησης βοηθούν την επεκτασιμότητα και
την αναπαραγωγιμότητα του εργαλείου.
\end{itemize}

Η περιγραφή βασίζεται στο πραγματικό codebase του έργου και όχι σε θεωρητικό
σχήμα ανεξάρτητο από την υλοποίηση
\cite{nlds_internal_md_docs,nlds_src_app_main}.

\section{Υπολογιστική στοίβα και αρχιτεκτονική}
\label{sec:ch2-stack-architecture}

Ο προσομοιωτής υλοποιείται σε Python και συνδυάζει εργαλεία για διεπαφή
χρήστη, αριθμητική ολοκλήρωση, συμβολικό χειρισμό εξισώσεων και παραγωγή
διαγραμμάτων. Η βασική στοίβα συνοψίζεται στον Πίνακα~\ref{tab:ch2-stack}.

\begin{table}[H]
\centering
\caption{Κύρια τεχνολογική στοίβα του προσομοιωτή}
\label{tab:ch2-stack}
\renewcommand{\arraystretch}{1.15}
\begin{tabular}{|l|p{0.64\textwidth}|}
\hline
\textbf{Συστατικό} & \textbf{Ρόλος στην υλοποίηση} \\
\hline
Python & Κεντρική γλώσσα υλοποίησης για διεπαφή, έλεγχο ροής, αριθμητική και exports. \\
\hline
Streamlit & Γραφικό περιβάλλον με tabs, widgets και \texttt{session\_state}. \\
\hline
NumPy & Πίνακες κατάστασης, διανυσματική αριθμητική και αριθμητική επεξεργασία δεδομένων. \\
\hline
SciPy (\texttt{solve\_ivp}) & Προσαρμοστικοί ολοκληρωτές RK45 και DOP853, καθώς και event detection. \\
\hline
SymPy & Parsing custom εξισώσεων, symbolic Jacobian και expression-based τομές Poincar\'e. \\
\hline
Numba & Επιτάχυνση fixed-step kernels, Lyapunov loops και sweep backends. \\
\hline
Pandas & Δομές δεδομένων για sweeps και εξαγωγές πινάκων. \\
\hline
Matplotlib & Διαγράμματα φάσης, χρονοσειρές, bifurcation και Lyapunov. \\
\hline
\end{tabular}
\end{table}

Η στοίβα αυτή φαίνεται άμεσα στα αρχεία εξαρτήσεων και στα κύρια modules της
εφαρμογής \cite{nlds_src_app_requirements,nlds_src_app_main}.

\begin{importantbox}
\textbf{Βασικός αρχιτεκτονικός διαχωρισμός}

\begin{itemize}
\item \textbf{Διεπαφή χρήστη:} \path{app/nlds_app.py} και
\path{app/ui/bifurcation_tab.py}. Εδώ ο χρήστης ορίζει το σύστημα, τους
χρόνους ολοκλήρωσης και τις ρυθμίσεις των διαγραμμάτων.
\item \textbf{Ορχήστρωση και cache:} \path{app/cache.py},
\path{app/logic/}, \path{app/sweep.py}, \path{app/export_utils.py}.
Εδώ γίνεται ο έλεγχος των ρυθμίσεων, η αποθήκευση ενδιάμεσων αποτελεσμάτων και
η επιλογή της κατάλληλης διαδρομής εκτέλεσης.
\item \textbf{Αριθμητικός πυρήνας:} \path{core/solver.py},
\path{core/symplectic_solver.py}, \path{core/poincare_sweep.py},
\path{core/lyapunov.py}, \path{core/numba_backend/}. Εδώ βρίσκονται οι
ολοκληρωτές και οι βασικοί αλγόριθμοι.
\end{itemize}
\end{importantbox}

Ο παραπάνω διαχωρισμός επιτρέπει να αλλάξει ο τρόπος εκτέλεσης ενός αριθμητικού
βρόχου χωρίς να χρειαστεί να αλλάξει η διεπαφή του χρήστη
\cite{nlds_src_app_main,nlds_src_ui_bifurcation_tab,nlds_src_app_cache,nlds_src_core_solver,nlds_src_core_lyapunov}.

\begin{figure}[H]
\centering
\fbox{\parbox[c][4.0cm][c]{0.88\textwidth}{\centering
Placeholder\\
Σχηματική ροή αρχιτεκτονικής:\\
sidebar / tabs $\rightarrow$ objects ρυθμίσεων $\rightarrow$ cache/orchestration
$\rightarrow$ numerical core}}
\caption{Placeholder για συνολικό διάγραμμα ροής της αρχιτεκτονικής του προσομοιωτή.}
\label{fig:ch2-architecture-placeholder}
\end{figure}

\section{Ροή χρήσης της εφαρμογής}
\label{sec:ch2-ui-flow}

\subsection{Ορισμός συστήματος και αρχικών δεδομένων}
\label{subsec:ch2-system-setup}

Η πρώτη επαφή του χρήστη με την εφαρμογή γίνεται από το sidebar. Εκεί
επιλέγεται αν θα χρησιμοποιηθεί ένα από τα έτοιμα συστήματα ή ένα custom
σύστημα, ποιες είναι οι αρχικές συνθήκες και ποιος solver θα χρησιμοποιηθεί.

\begin{guidebox}
Στην αρχική ρύθμιση του προβλήματος ο χρήστης δίνει:
\begin{enumerate}
\item το είδος του συστήματος,
\item τις παραμέτρους του μοντέλου,
\item τις αρχικές συνθήκες,
\item και, στην custom περίπτωση, τα ονόματα μεταβλητών και τις εξισώσεις.
\end{enumerate}
\end{guidebox}

Για custom συστήματα η εφαρμογή επιτρέπει και symbolic παραγωγή Ιακωβιανού,
ώστε ο ίδιος ο χρήστης να επιλέξει αργότερα αν θα χρησιμοποιηθεί αναλυτικός ή
αριθμητικός Jacobian στους υπολογισμούς Lyapunov
\cite{nlds_src_app_main,nlds_src_app_numba_custom,nlds_src_core_lyapunov}.

\begin{figure}[H]
\centering
\begin{subfigure}{0.32\textwidth}
  \centering
  \includegraphics[width=\textwidth]{figures/ch2/s1_system_selection.png}
  \caption{Επιλογή συστήματος και solver kind.}
  \label{fig:ch2-s1-system}
\end{subfigure}
\hfill
\begin{subfigure}{0.32\textwidth}
  \centering
  \includegraphics[width=\textwidth]{figures/ch2/s2_initial_and_custom_definitions.png}
  \caption{Αρχικές συνθήκες και custom εξισώσεις.}
  \label{fig:ch2-s2-custom-def}
\end{subfigure}
\hfill
\begin{subfigure}{0.32\textwidth}
  \centering
  \includegraphics[width=\textwidth]{figures/ch2/s3_symbolic_jacobian.png}
  \caption{Προβολή symbolic Jacobian.}
  \label{fig:ch2-s3-jacobian}
\end{subfigure}
\caption{Βασικά βήματα ορισμού του συστήματος στο sidebar και στα πρώτα blocks της εφαρμογής.}
\label{fig:ch2-tab-setup}
\end{figure}

\subsection{Ολοκλήρωση τροχιάς και οπτικοποίηση}
\label{subsec:ch2-integration-visualization}

Αφού οριστεί το πρόβλημα, ο χρήστης επιλέγει το χρονικό παράθυρο ολοκλήρωσης,
το βήμα \(\Delta t\), τους άξονες προβολής και τις ανοχές του solver όπου αυτό
χρειάζεται. Η βασική χρονική δειγματοληψία γράφεται ως
\begin{equation}
t_k=t_0+k\Delta t,
\qquad k=0,1,\dots,N.
\end{equation}
Με αυτά τα δεδομένα παράγονται η τροχιά αναφοράς, το φασικό πορτρέτο και οι
χρονοσειρές.

\begin{guidebox}
Στο πρώτο tab συγκεντρώνονται οι ρυθμίσεις που αφορούν:
\begin{enumerate}
\item τον χρόνο ολοκλήρωσης,
\item τις παραμέτρους του συστήματος,
\item την αφαίρεση του μεταβατικού μέρους,
\item την προβολή σε 2D ή 3D άξονες,
\item και την αποθήκευση των ρυθμίσεων του run.
\end{enumerate}
\end{guidebox}

\begin{figure}[H]
\centering
\begin{subfigure}{0.32\textwidth}
  \centering
  \includegraphics[width=\textwidth]{figures/ch2/s4_integration_controls.png}
  \caption{Χρονικό παράθυρο ολοκλήρωσης.}
  \label{fig:ch2-s4-integration}
\end{subfigure}
\hfill
\begin{subfigure}{0.32\textwidth}
  \centering
  \includegraphics[width=\textwidth]{figures/ch2/s5_system_parameters_custom.png}
  \caption{Παράμετροι συστήματος.}
  \label{fig:ch2-s5-params}
\end{subfigure}
\hfill
\begin{subfigure}{0.32\textwidth}
  \centering
  \includegraphics[width=\textwidth]{figures/ch2/s6_plot_settings_lyapunov.png}
  \caption{Plot και Lyapunov settings.}
  \label{fig:ch2-s6-plot}
\end{subfigure}
\caption{Κύριες ρυθμίσεις ολοκλήρωσης και post-processing στο πρώτο tab.}
\label{fig:ch2-tab1-top}
\end{figure}

\begin{figure}[H]
\centering
\begin{subfigure}{0.32\textwidth}
  \centering
  \includegraphics[width=\textwidth]{figures/ch2/s7_axis_selection.png}
  \caption{Επιλογή αξόνων.}
  \label{fig:ch2-s7-axis}
\end{subfigure}
\hfill
\begin{subfigure}{0.32\textwidth}
  \centering
  \includegraphics[width=\textwidth]{figures/ch2/s8_solver_tolerances.png}
  \caption{Ανοχές solver.}
  \label{fig:ch2-s8-tols}
\end{subfigure}
\hfill
\begin{subfigure}{0.32\textwidth}
  \centering
  \includegraphics[width=\textwidth]{figures/ch2/s9_save_configuration.png}
  \caption{Αποθήκευση configuration.}
  \label{fig:ch2-s9-save}
\end{subfigure}
\caption{Ρυθμίσεις προβολής, ανοχών και αποθήκευσης του run.}
\label{fig:ch2-tab1-bottom}
\end{figure}

Το δεύτερο tab δεν λύνει ξανά το πρόβλημα από την αρχή. Πατά πάνω στην ήδη
υπολογισμένη τροχιά και επιτρέπει να επιλεγεί ένα μικρότερο χρονικό παράθυρο
για καθαρότερη απεικόνιση των χρονοσειρών
\cite{nlds_src_app_main}.

\begin{figure}[H]
\centering
\includegraphics[width=0.56\textwidth]{figures/ch2/s10_time_series_window.png}
\caption{Το δεύτερο tab επιτρέπει επιλογή χρονικού παραθύρου πάνω στην ήδη
υπολογισμένη χρονοσειρά.}
\label{fig:ch2-s10-ts}
\end{figure}

\subsection{Παραμετρικές σαρώσεις}
\label{subsec:ch2-sweeps-ui}

Στο τρίτο tab ο χρήστης δεν ακολουθεί μία μόνο τροχιά, αλλά επαναλαμβάνει την
ολοκλήρωση για πολλές τιμές μιας παραμέτρου. Αν η παράμετρος γραφεί ως
\(\mu\), τότε τα σημεία της σάρωσης δίνονται από τη σχέση
\begin{equation}
\mu_i=\mu_{\min}+i\,\Delta\mu.
\end{equation}
Για κάθε τιμή της παραμέτρου παράγεται είτε σύνολο σημείων για διάγραμμα
διακλάδωσης είτε ένα φάσμα εκθετών Lyapunov.

\begin{guidebox}
Το τρίτο tab χωρίζεται λειτουργικά σε δύο μέρη:
\begin{enumerate}
\item ρυθμίσεις για bifurcation sweep με τομή Poincar\'e ή local extrema,
\item ρυθμίσεις για Lyapunov sweep με QR interval και κανόνες clipping.
\end{enumerate}
\end{guidebox}

\begin{figure}[H]
\centering
\begin{subfigure}{0.32\textwidth}
  \centering
  \includegraphics[width=\textwidth]{figures/ch2/s11_tab3_top_setup.png}
  \caption{Γενικές ρυθμίσεις sweep.}
  \label{fig:ch2-s11-tab3-top}
\end{subfigure}
\hfill
\begin{subfigure}{0.32\textwidth}
  \centering
  \includegraphics[width=\textwidth]{figures/ch2/s12_bifurcation_parallel_section.png}
  \caption{Ρυθμίσεις bifurcation sweep.}
  \label{fig:ch2-s12-bif-section}
\end{subfigure}
\hfill
\begin{subfigure}{0.32\textwidth}
  \centering
  \includegraphics[width=\textwidth]{figures/ch2/s13_bifurcation_method_earlystop.png}
  \caption{Method και early stop.}
  \label{fig:ch2-s13-bif-method}
\end{subfigure}
\caption{Κύριες επιλογές για bifurcation sweep στο τρίτο tab.}
\label{fig:ch2-tab3-top}
\end{figure}

\begin{figure}[H]
\centering
\begin{subfigure}{0.48\textwidth}
  \centering
  \includegraphics[width=\textwidth]{figures/ch2/s14_lyapunov_sweep_settings.png}
  \caption{Ρυθμίσεις Lyapunov sweep.}
  \label{fig:ch2-s14-lya-settings}
\end{subfigure}
\hfill
\begin{subfigure}{0.48\textwidth}
  \centering
  \includegraphics[width=\textwidth]{figures/ch2/s15_lyapunov_clip_and_run.png}
  \caption{Clipping και εκτέλεση sweep.}
  \label{fig:ch2-s15-lya-run}
\end{subfigure}
\caption{Ρυθμίσεις Lyapunov sweep στο τρίτο tab.}
\label{fig:ch2-tab3-bottom}
\end{figure}

\section{Από τη διεπαφή στον αριθμητικό πυρήνα}
\label{sec:ch2-ui-to-core}

\begin{definitionbox}{Objects ρυθμίσεων του run}{nlds_src_app_params}
Οι βασικές ρυθμίσεις ενός run συγκεντρώνονται σε μικρά dataclasses. Τα πιο
κεντρικά είναι τα \texttt{SystemConfig}, \texttt{InitialConditions},
\texttt{IntegrationConfig}, \texttt{SolverTolerances},
\texttt{SweepRunConfig} και \texttt{LyapunovConfig}. Με αυτόν τον τρόπο η
εφαρμογή μεταφέρει οργανωμένα τα δεδομένα από τη διεπαφή στον πυρήνα.
\end{definitionbox}

Στην πράξη, ο χρήστης δεν επικοινωνεί απευθείας με τις numerical routines.
Πρώτα δημιουργούνται τα objects των ρυθμίσεων και έπειτα επιλέγεται το σωστό
μονοπάτι εκτέλεσης.

\begin{importantbox}
\textbf{Οι τρεις βασικές εκτελεστικές διαδρομές}

\begin{itemize}
\item \textbf{Μοναδική τροχιά:}
\path{app/nlds_app.py}
\(\rightarrow\)
\path{app/cache.py:solve_cached}
\(\rightarrow\)
\path{core/solver.py} ή \path{core/symplectic_solver.py}.
\item \textbf{Μεμονωμένος υπολογισμός Lyapunov:}
\path{app/nlds_app.py}
\(\rightarrow\)
\path{app/logic/lyapunov_cached.py}
\(\rightarrow\)
\path{core/lyapunov.py}.
\item \textbf{Sweep παραμέτρων:}
\path{app/ui/bifurcation_tab.py}
\(\rightarrow\)
\path{app/sweep.py} ή \path{app/logic/lyapunov_sweep.py}
\(\rightarrow\)
\path{core/poincare_sweep.py} ή \path{core/lyapunov.py}.
\end{itemize}
\end{importantbox}

\subsection{Μοναδική τροχιά}
\label{subsec:ch2-single-run-path}

Για ένα απλό run, η διεπαφή συγκεντρώνει τα στοιχεία του προβλήματος και
καλεί τη συνάρτηση \texttt{solve\_cached}. Η συνάρτηση αυτή αποφασίζει αν θα
χρησιμοποιηθεί adaptive solver, fixed-step RK4 ή symplectic ολοκληρωτής,
ανάλογα με το \texttt{solver\_kind} και το είδος του συστήματος
\cite{nlds_src_app_cache,nlds_src_core_solver,nlds_src_core_symplectic}.

Στο σημείο αυτό εμφανίζεται και το πρώτο απλό μαθηματικό φίλτρο της
υλοποίησης: σε χαμιλτονιανά runs το state πρέπει να μπορεί να γραφεί ως
\([\mathbf{q},\mathbf{p}]\), δηλαδή ως ζεύγος γενικευμένων θέσεων και ορμών.
Γι' αυτό ο symplectic solver ενεργοποιείται μόνο όταν ο αριθμός μεταβλητών
είναι άρτιος και το σύστημα έχει την κατάλληλη μορφή.

\subsection{Υπολογισμός φάσματος Lyapunov}
\label{subsec:ch2-lyapunov-path}

Για τους εκθέτες Lyapunov η διεπαφή χρησιμοποιεί διαφορετικό μονοπάτι. Το
\path{app/logic/lyapunov_cached.py} προετοιμάζει το σύστημα, επιλέγει αν θα
χρησιμοποιηθεί analytic ή finite-difference Jacobian και τελικά καλεί το
\path{core/lyapunov.py} \cite{nlds_src_logic_lyapunov_cached,nlds_src_core_lyapunov}.

Η πιο βασική χρονική ρύθμιση εδώ είναι το QR interval. Αν \(\Delta t\) είναι
το αριθμητικό βήμα και \(q\) το πλήθος βημάτων ανά ορθοκανονικοποίηση, τότε
\begin{equation}
\Delta T_{\mathrm{QR}} = q\,\Delta t.
\end{equation}
Η σχέση αυτή χρησιμοποιείται άμεσα από τον κώδικα για να περάσουμε από τη
ρύθμιση της διεπαφής στο πραγματικό αριθμητικό χρονοδιάστημα του υπολογισμού.

\subsection{Sweeps παραμέτρων}
\label{subsec:ch2-sweep-path}

Στα sweeps το ίδιο πείραμα επαναλαμβάνεται πολλές φορές. Στο bifurcation μέρος
του τρίτου tab, η \texttt{run\_sweep\_chunk} συλλέγει είτε τομές Poincar\'e
είτε local extrema, ανάλογα με τη ρύθμιση που επέλεξε ο χρήστης
\cite{nlds_src_ui_bifurcation_tab,nlds_src_app_sweep,nlds_src_core_poincare}.

Στο Lyapunov sweep, η ροή είναι αντίστοιχη αλλά ο τελικός στόχος δεν είναι ένα
σύνολο σημείων τομής. Για κάθε τιμή της παραμέτρου υπολογίζεται ένα φάσμα
εκθετών και αποθηκεύεται συνήθως ο μεγαλύτερος εκθέτης ή το πλήρες διάνυσμα,
ανάλογα με το mode της εκτέλεσης
\cite{nlds_src_logic_lyapunov_sweep,nlds_src_core_lyapunov}.

\section{Νεότερα στοιχεία της τρέχουσας υλοποίησης}
\label{sec:ch2-current-implementation}

Το σημερινό codebase έχει αρκετά στοιχεία που δεν φαίνονται αμέσως από το
θεωρητικό κεφάλαιο, αλλά είναι σημαντικά για την πρακτική λειτουργία της
εφαρμογής.

\subsection{Φραγμένη αποθήκευση τροχιάς}
\label{subsec:ch2-bounded-trajectory}

Σε ένα μεγάλο run δεν είναι πάντα χρήσιμο να αποθηκεύονται όλα τα σημεία της
τροχιάς. Για τον λόγο αυτό το \texttt{IntegrationConfig} περιλαμβάνει την
παράμετρο \texttt{max\_store\_steps}. Έτσι, ο solver μπορεί να ολοκληρώνει σε
μεγάλο χρονικό διάστημα, αλλά η μνήμη που δεσμεύεται για το τελικό αποθηκευμένο
αποτέλεσμα να παραμένει ελεγχόμενη
\cite{nlds_src_app_params,nlds_src_app_cache,nlds_src_core_solver,nlds_src_core_symplectic}.

Η ίδια λογική συνεχίζεται και στην οπτικοποίηση, όπου το plotting
χρησιμοποιεί decimation ώστε να μη στέλνονται στο γράφημα πολύ περισσότερα
σημεία από όσα μπορούν να διαβαστούν εύκολα.

\subsection{Μεγάλα sweeps με bounded memory}
\label{subsec:ch2-bounded-sweeps}

Στα μεγάλα bifurcation sweeps, η εφαρμογή δεν κρατά απεριόριστα όλες τις
παλιές γραμμές στη RAM. Διατηρεί ένα πρόσφατο κομμάτι δεδομένων σε πλήρη
ανάλυση και ένα reservoir sample από παλαιότερα σημεία, ώστε το διάγραμμα να
μένει αντιπροσωπευτικό χωρίς να μεγαλώνει απεριόριστα η μνήμη
\cite{nlds_src_ui_bifurcation_tab,nlds_src_logic_reservoir}.

Η λεπτομερής ανάλυση μνήμης για αυτά τα σημεία δίνεται στο
Παράρτημα~\ref{app:ch2-memory}.

\subsection{Αποθήκευση ρυθμίσεων και αναπαραγωγιμότητα}
\label{subsec:ch2-reproducibility}

Η εφαρμογή επιτρέπει αποθήκευση τόσο ενός απλού run όσο και ενός sweep σε
μορφή JSON. Το \texttt{app/export\_utils.py} κατασκευάζει οργανωμένα blocks με
πληροφορίες για το σύστημα, τις αρχικές συνθήκες, τον solver, τα post-process
settings και, όταν είναι διαθέσιμο, το commit του repository
\cite{nlds_src_app_export}.

Αυτό είναι χρήσιμο γιατί το αποτέλεσμα ενός πειράματος δεν εξαρτάται μόνο από
τις εξισώσεις. Εξαρτάται επίσης από:
\begin{itemize}
\item το \(\Delta t\),
\item το transient cut,
\item τις ανοχές του solver,
\item το QR interval,
\item και το mode του sweep.
\end{itemize}

Με την αποθήκευση αυτών των ρυθμίσεων το πείραμα μπορεί να επαναληφθεί πολύ πιο
αξιόπιστα.

\subsection{Compiled paths και ασφαλή fallbacks}
\label{subsec:ch2-numba-short}

Η εφαρμογή δεν τρέχει όλον τον κώδικα σε compiled μορφή. Η διεπαφή, ο έλεγχος
ροής και η παραγωγή γραφημάτων παραμένουν σε Python. Εκεί που χρειάζεται
επιτάχυνση, ενεργοποιούνται Numba kernels για fixed-step RK4, symplectic runs,
Lyapunov υπολογισμούς και ορισμένα sweep backends
\cite{nlds_src_core_numba,nlds_src_app_numba_custom}.

Αν το compiled μονοπάτι δεν είναι διαθέσιμο ή αν μια custom περίπτωση δεν
μπορεί να μεταγλωττιστεί με ασφάλεια, ο κώδικας επιστρέφει στην κανονική
Python διαδρομή. Με αυτόν τον τρόπο η εφαρμογή παραμένει λειτουργική χωρίς να
θυσιάζει τη γενικότητα των custom ορισμών.

Η πιο τεχνική περιγραφή των compiled paths δίνεται στο
Παράρτημα~\ref{app:ch2-numba}.

\section{Συνολική εικόνα του κεφαλαίου}
\label{sec:ch2-closing}

Το βασικό συμπέρασμα του κεφαλαίου είναι ότι ο προσομοιωτής δεν είναι απλώς ένα
σύνολο μεμονωμένων scripts. Είναι μια ενιαία ροή εργασίας στην οποία:
\begin{itemize}
\item η διεπαφή συγκεντρώνει οργανωμένα τις ρυθμίσεις,
\item το layer ορχήστρωσης επιλέγει το κατάλληλο μονοπάτι εκτέλεσης,
\item και ο αριθμητικός πυρήνας εκτελεί το ουσιαστικό υπολογιστικό έργο.
\end{itemize}

Η κύρια ροή κρατήθηκε εδώ όσο γίνεται απλή. Οι αναλυτικότεροι τύποι
πολυπλοκότητας, τα όρια μνήμης και η τεχνική περιγραφή της Numba/JIT
υλοποίησης μεταφέρονται στο επόμενο παράρτημα.
